% 共享章节数据 — 由 scripts/filter_config.json + notes_config.json 生成
% 各布局（print/reader）自行定义 \slidesdata 和 \notedata 来解读


% ==== 引言（9 页）====
\slidesdata{1引言}{5,39,40,41,42,43,44,45,46}

% ==== 数学基础（13 页）====
\slidesdata{2数学基础}{3,4,5,6,7,8,9,10,11,12,13,14,15}

% ==== 凸分析（28 页）====
\slidesdata{3凸分析}{3,4,5,8,9,11,12,14,15,17,18,21,23,24,25,27,29,30,34,39,40,41,42,43,44,45,46,47}

% ==== 线性规划基本性质（26 页）====
\slidesdata{4线性规划基本性质}{3,4,5,6,7,8,9,10,11,12,13,14,15,16,17,21,22,23,24,25,26,27,28,29,30,31}

% ==== 单纯形方法（40 页）====
\slidesdata{5单纯形方法}{2,3,4,5,6,7,12,15,16,24,25,26,27,28,29,30,35,36,37,38,39,40,41,42,43,44,45,46,47,48,49,50,51,52,53,54,55,56,57,58}

% ==== 对偶理论（58 页）====
\slidesdata{6对偶理论}{3,4,5,6,7,8,9,10,11,15,16,17,18,19,21,22,23,24,25,26,27,28,29,30,31,32,33,40,41,42,43,44,45,46,47,48,49,50,51,52,53,54,55,56,57,58,59,60,61,62,63,64,65,66,67,68,71,81}

% ==== 最优性条件（43 页）====
\slidesdata{7最优性条件}{3,4,5,7,13,14,15,16,17,18,19,20,21,22,23,24,25,26,27,28,29,34,35,36,37,38,39,40,41,42,43,48,49,50,51,52,53,54,55,56,57,58,59}
\notedata{7最优性条件}{\section*{FJ条件 vs KKT条件 对比}

\textbf{FJ条件}：存在不全为零的 $\lambda_0,\dots,\lambda_m,\mu_1,\dots,\mu_l$ 使得
$\lambda_0\nabla f + \sum\lambda_i\nabla g_i + \sum\mu_j\nabla h_j = 0$，
$\lambda_i \ge 0$，$\lambda_i g_i = 0$。
注意：$\lambda_0$ 可以为 0，此时不含目标函数信息。

\textbf{KKT条件}：增加约束规格后，保证 $\lambda_0 \neq 0$（可归一化为1）。
$\nabla f + \sum\lambda_i\nabla g_i + \sum\mu_j\nabla h_j = 0$，
$\lambda_i \ge 0$，$\lambda_i g_i = 0$。
即 Lagrange 函数 $L(x,\lambda,\mu)$ 的驻点条件。

\textbf{核心区别：}FJ中$\lambda_0$可能为0（不含目标函数），KKT通过约束规格保证$\lambda_0 \neq 0$。

真题链接：2025 Q6 — 先求所有FJ点，再证不存在KKT点。}

% ==== 算法（19 页）====
\slidesdata{8算法}{3,4,5,6,7,8,9,10,14,15,16,17,18,19,20,21,22,23,24}

% ==== 一维搜索（17 页）====
\slidesdata{9一维搜索}{3,4,5,6,7,26,27,51,52,53,54,55,56,57,58,59,60}

% ==== 使用导数的优化算法（79 页）====
\slidesdata{10使用导数的优化算法}{40,41,42,43,44,45,46,47,48,49,50,51,52,53,54,55,56,57,58,59,60,61,62,63,64,65,66,67,68,69,70,71,72,73,74,75,76,77,78,79,80,81,82,83,84,85,86,87,88,89,90,91,92,93,94,95,96,97,98,99,100,101,102,103,104,105,106,107,108,109,110,111,112,113,114,115,116,117,118}
\notedata{10使用导数的优化算法}{\section*{拟牛顿法公式速查}

\textbf{DFP法（对Hessian逆矩阵$H$的秩2校正）}
\begin{align*}
  s_k &= x_{k+1} - x_k,\qquad y_k = g_{k+1} - g_k \\
  H_{k+1} &= H_k + \frac{s_k s_k^T}{s_k^T y_k}
            - \frac{H_k y_k y_k^T H_k}{y_k^T H_k y_k}
\end{align*}
搜索方向: $d_k = -H_k \nabla f(x_k)$

\textbf{BFGS法（对Hessian矩阵$B$的秩2校正）}
\begin{align*}
  B_{k+1} &= B_k + \frac{y_k y_k^T}{y_k^T s_k}
            - \frac{B_k s_k s_k^T B_k}{s_k^T B_k s_k}
\end{align*}
搜索方向: 解$B_k d_k = -\nabla f(x_k)$

\textbf{关键区别：} DFP校$H$（Hessian逆），BFGS校$B$（Hessian本身）。

真题链接：2025 Q4 — DFP法完整两轮迭代。}

% ==== 无约束优化直接方法（0 页）====
\slidesdata{11无约束优化直接方法}{}

% ==== 可行方向法（60 页）====
\slidesdata{12可行方向法}{2,3,4,5,6,7,8,9,10,11,12,13,14,15,16,17,18,19,20,21,22,23,24,25,26,27,28,29,30,31,32,33,39,40,41,42,43,44,45,46,47,48,49,50,51,52,53,54,55,56,57,58,59,60,61,62,63,64,65,66}
\notedata{12可行方向法}{\section*{Rosen梯度投影法 算法步骤}

投影矩阵: $P = I - M^T(MM^T)^{-1}M$
搜索方向: $d = -P\,\nabla f(x)$

\textbf{算法流程}
\begin{enumerate}
  \item 在当前点$x_k$识别\textbf{积极约束}（取等号的约束），构造积极约束系数矩阵$M$
  \item 计算投影矩阵 $P = I - M^T(MM^T)^{-1}M$
  \item 计算搜索方向 $d_k = -P \nabla f(x_k)$
  \item \textbf{若$d_k \neq 0$}：沿$d_k$做一维搜索，求步长$\lambda_k$，$x_{k+1} = x_k + \lambda_k d_k$，回到1
  \item \textbf{若$d_k = 0$}：计算$w = -(MM^T)^{-1}M\nabla f(x_k)$
        \begin{itemize}
          \item 若$w$的所有分量$\ge 0$ → $x_k$为\textbf{KKT点}，停止
          \item 若$w$有负分量 → 去掉对应的积极约束行，重新计算$P$
        \end{itemize}
\end{enumerate}

\textbf{Zoutendijk法（补充）}
解LP子问题求下降可行方向：
\[ \min\ \nabla f(x)^T d,\ \text{s.t.}\ \nabla g_i(x)^T d + g_i(x) \le 0,\ -1 \le d_j \le 1 \]
若最优值$<0$得下降可行方向；若$=0$则为FJ/KKT点。

真题链接：2025 Q5 — Rosen梯度投影法求解约束优化。}

% ==== 罚函数法（48 页）====
\slidesdata{13罚函数法}{2,3,4,5,6,7,8,9,10,11,12,13,14,15,16,17,18,19,20,21,22,23,24,25,26,27,28,29,30,31,32,33,34,35,36,37,38,39,40,41,42,43,44,45,46,47,48,49}
